\documentclass{article}
\usepackage{comment}
\usepackage{amsmath}
\usepackage{amsfonts}
\usepackage{sectsty}
\usepackage{hyperref}
\usepackage{fancyhdr}
\usepackage[top=1in,bottom=1in,left=1in,right=1in]{geometry}
\usepackage{float}
\usepackage{graphicx}
\usepackage{tabularx}
\usepackage{mathtools}

\title{Assignment 3}
\author{Matthew Farah, \href{mailto:farahm11@mcmaster.ca}{$\langle$farahm11@mcmaster.ca$\rangle$}, 400468251}
\date{\today}

\hypersetup{
        colorlinks=true,
        linkcolor=blue,
        filecolor=magenta,
        urlcolor=blue,
        citecolor=blue,
	anchorcolor=blue
}

\fancyhead{}
\fancyhead[L]{Matthew Farah, \href{mailto:farahm11@mcmaster.ca}{$\langle$farahm11@mcmaster.ca$\rangle$}, 400468251}
\fancyhead[R]{Assignment 3}

\newenvironment{directsubsubs}{
  \makeatletter
  \renewcommand\thesubsubsection{\thesection.\Roman{subsubsection}}
  \makeatother
}{
  \makeatletter
  \renewcommand\thesubsubsection{\thesubsection.\Roman{subsubsection}}
  \makeatother
}

\renewcommand{\thesubsection}{\thesection.\alph{subsection}}
\renewcommand{\thesubsubsection}{\thesection.\alph{subsection}.\Roman{subsubsection}}
\makeatletter
\DeclareFontFamily{U}{MnSymbolA}{}
\DeclareFontShape{U}{MnSymbolA}{m}{n}{
    <-6>  MnSymbolA5
   <6-7>  MnSymbolA6
   <7-8>  MnSymbolA7
   <8-9>  MnSymbolA8
   <9-10> MnSymbolA9
  <10-12> MnSymbolA10
  <12->   MnSymbolA12}{}
\DeclareFontShape{U}{MnSymbolA}{b}{n}{
    <-6>  MnSymbolA-Bold5
   <6-7>  MnSymbolA-Bold6
   <7-8>  MnSymbolA-Bold7
<8-9>  MnSymbolA-Bold8
   <9-10> MnSymbolA-Bold9
  <10-12> MnSymbolA-Bold10
  <12->   MnSymbolA-Bold12}{}
\DeclareSymbolFont{MnSyA}{U}{MnSymbolA}{m}{n}
\SetSymbolFont{MnSyA}{bold}{U}{MnSymbolA}{b}{n}

\DeclareRobustCommand{\overleftharpoon}{\mathpalette{\overarrow@\leftharpoonfill@}}
\DeclareRobustCommand{\overrightharpoon}{\mathpalette{\overarrow@\rightharpoonfill@}}
\def\leftharpoonfill@{\arrowfill@\leftharpoondown\mn@relbar\mn@relbar}
\def\rightharpoonfill@{\arrowfill@\mn@relbar\mn@relbar\rightharpoonup}

\DeclareMathSymbol{\leftharpoondown}{\mathrel}{MnSyA}{'112}
\DeclareMathSymbol{\rightharpoonup}{\mathrel}{MnSyA}{'100}
\DeclareMathSymbol{\mn@relbar}{\mathrel}{MnSyA}{'320}
\makeatother


\newcolumntype{Y}{>{\centering\arraybackslash}X}
\renewcommand{\arraystretch}{1.8}

\sectionfont{\fontsize{12}{12}\bfseries}
\subsectionfont{\fontsize{10}{12}\normalfont}
\subsubsectionfont{\fontsize{10}{12}\normalfont}

\newcommand{\harpoon}[1]{\overrightharpoonup{#1}}

\begin{document}
\maketitle
\pagestyle{fancy}

\begin{center}
	\Large{\underline{System Representations}}
\end{center}

\section[Question ~\thesection]{Given the impulse response $h(n) = \delta(n) + \frac{1}{2}\delta(n - 1) + \frac{1}{4}\delta(n - 2)$}

\subsection[~\thesubsection]{Determine the corresponding difference equation}

Since we are not guarenteed that the system is FIR, we must recover $y(n)$ from the impulse response using the convolutional equation.

\begin{align*}
	y(n) &= \sum_{k = -\infty}^{\infty}(h(n - k)x(k)) \\
	&= \sum_{k = -\infty}^{\infty}((\delta(n - k) + \frac{1}{2}\delta(n - k - 1) + \frac{1}{4}\delta(n - k - 2))x(n)) \\
	&= \sum_{k = -\infty}^{\infty}(\delta(n - k)x(n) + \frac{1}{2}\delta(n - k - 1)x(n) + \frac{1}{4}\delta(n - k - 2)x(n)) \\
	&= \delta(0)x(n) + \frac{1}{2}\delta(0)x(n - 1) + \frac{1}{4}\delta(0)x(n - 2) \\
	&= x(n) + \frac{1}{2}x(n - 1) + \frac{1}{4}x(n - 2)
\end{align*}

Thus, the system's difference equation is given by $y(n) = x(n) + \frac{1}{2}x(n - 1) + \frac{1}{4}x(n - 2)$.

\subsection[~\thesubsection]{Determine the $(A, B, C, D)$ representation of the system}

\subsubsection[~\thesubsubsection]{Find $\overrightharpoon{S(n)}$}

\begin{align*}
	&\text{Let:}& \text{Then:} \\
	&s_1(n) = x(n - 2) & s_1(n + 1) = x(n - 1) = s_2(n) \\
	&s_2(n) = x(n - 1) & s_2(n + 1) = x(n)
\end{align*}

Thus, let $\overrightharpoon{S(n + 1)} = \begin{psmallmatrix} s_1(n) \\ s_2(n) \end{psmallmatrix}$

\subsubsection[~\thesubsubsection]{Find $\overrightharpoon{S(n + 1)}$}

\begin{align*}
	\overrightharpoon{S(n + 1)} &= \begin{pmatrix} s_1(n + 1) \\ s_2(n + 1) \end{pmatrix} \\
	&= \begin{pmatrix} s_2(n) \\ x(n) \end{pmatrix} \\
	&= \begin{pmatrix} 0 & 1 \\ 0 & 0 \end{pmatrix}\overrightharpoon{S(n)} + \begin{pmatrix} 0 \\ 1 \end{pmatrix}x(n)
\end{align*}

Therefore, our $\boldsymbol{A}$ and $\boldsymbol{B}$ matrices are $\begin{psmallmatrix} 0 & 1 \\ 0 & 0 \end{psmallmatrix}$ and $\begin{psmallmatrix} 0 \\ 1 \end{psmallmatrix}$ respectively.

\subsubsection[~\thesubsubsection]{Find $y(n)$}

\begin{align*}
	y(n) &= \frac{1}{4}x(n - 2) + \frac{1}{2}x(n - 1) + x(n) \\
	&= \frac{1}{4}s_1(n) + \frac{1}{2}s_2(n) + x(n) \\
	&= \begin{pmatrix} \frac{1}{4} & \frac{1}{2} \end{pmatrix}\overrightharpoon{S(n)} + \begin{pmatrix} 1 \end{pmatrix}x(n)
\end{align*}

Therefore, our $\boldsymbol{C}$ and $\boldsymbol{D}$ matrices are $\begin{psmallmatrix} \frac{1}{4} & \frac{1}{2} \end{psmallmatrix}$ and $\begin{psmallmatrix} 1 \end{psmallmatrix}$ respectively.

Putting it all together, our system can be represented in state space form as:

\begin{align*}
	\overrightharpoon{S(n + 1)} &= \begin{pmatrix} 0 & 1 \\ 0 & 0 \end{pmatrix} + \begin{pmatrix} 0 \\ 1 \end{pmatrix} \\\\
	y(n) &= \begin{pmatrix} \frac{1}{4} & \frac{1}{2} \end{pmatrix}\overrightharpoon{S(n)} + \begin{pmatrix} 1 \end{pmatrix}x(n)
\end{align*}

\subsection[~\thesubsection]{Compute the output of the system for the input $x(n) = \delta(n) + \delta(n - 2)$:}

\subsubsection[~\thesubsubsection]{Using the difference equation}

\begin{align*}
	y(n) &= x(n) + \frac{1}{2}x(n - 1) + \frac{1}{4}x(n - 2) \\
	&= \delta(n) + \delta(n - 2) + \frac{1}{2}(\delta(n - 1) + \delta(n - 3)) + \frac{1}{4}(\delta(n - 2) + \delta(n - 4)) \\
	&= \delta(n) + \frac{1}{2}\delta(n - 1) + \frac{5}{4}\delta(n - 2) + \frac{1}{2}\delta(n - 3) + \frac{1}{4}\delta(n - 4) \\
\end{align*}

Thus, our system's outputs $y(n)$ for the input $x(n) = \delta(n) + \delta(n - 2)$ is given by:

\begin{align*}
	y(n) &= \delta(n) + \frac{1}{2}\delta(n - 1) + \frac{5}{4}\delta(n - 2) + \frac{1}{2}\delta(n - 3) + \frac{1}{4}\delta(n - 4) \\\\
	\intertext{\centering Or}
	y(n) &= \begin{cases} 1, \; \text{if} \; n = 0 \\ \frac{1}{2} \; \text{if} \; n = 1 \\ \frac{5}{4} \; \text{if} \; n = 2 \\ \frac{1}{2} \; \text{if} \; n = 3 \\ \frac{1}{4} \; \text{if} \; n = 4 \\0, \; \text{otherwise} \end{cases}
\end{align*}

\subsubsection[~\thesubsubsection]{Using the state-space equations}

\begin{align*}
	&\underline{\text{Current State} \; (\overrightharpoon{S(n)})}& &\underline{\text{Current Output} \; (y(n))} \\ \\
	&\overrightharpoon{S(0)} = \overrightharpoon{0}& &y(0) = \begin{pmatrix} \frac{1}{4} & \frac{1}{2} \end{pmatrix}\overrightharpoon{S(0)} + \begin{pmatrix} 1 \end{pmatrix}x(0) = 1 \\
	&\overrightharpoon{S(1)} = \begin{pmatrix} 0 & 1 \\ 0 & 0 \end{pmatrix}\overrightharpoon{S(0)} + \begin{pmatrix} 0 \\ 1 \end{pmatrix}x(0) = \begin{pmatrix} 0 \\ 1 \end{pmatrix}& &y(1) = \begin{pmatrix} \frac{1}{4} & \frac{1}{2} \end{pmatrix}\overrightharpoon{S(1)} + \begin{pmatrix} 1 \end{pmatrix}x(1) = \frac{1}{2} \\
	&\overrightharpoon{S(2)} = \begin{pmatrix} 0 & 1 \\ 0 & 0 \end{pmatrix}\overrightharpoon{S(1)} + \begin{pmatrix} 0 \\ 1 \end{pmatrix}x(1) = \begin{pmatrix} 1 \\ 0 \end{pmatrix}& &y(2) = \begin{pmatrix} \frac{1}{4} & \frac{1}{2} \end{pmatrix}\overrightharpoon{S(2)} + \begin{pmatrix} 1 \end{pmatrix}x(2) = \frac{5}{4} \\
	&\overrightharpoon{S(3)} = \begin{pmatrix} 0 & 1 \\ 0 & 0 \end{pmatrix}\overrightharpoon{S(2)} + \begin{pmatrix} 0 \\ 1 \end{pmatrix}x(2) = \begin{pmatrix} 0 \\ 1 \end{pmatrix}& &y(3) = \begin{pmatrix} \frac{1}{4} & \frac{1}{2} \end{pmatrix}\overrightharpoon{S(3)} + \begin{pmatrix} 1 \end{pmatrix}x(3) = \frac{1}{2} \\
	&\overrightharpoon{S(4)} = \begin{pmatrix} 0 & 1 \\ 0 & 0 \end{pmatrix}\overrightharpoon{S(3)} + \begin{pmatrix} 0 \\ 1 \end{pmatrix}x(3) = \begin{pmatrix} 1 \\ 0 \end{pmatrix}& &y(4) = \begin{pmatrix} \frac{1}{4} & \frac{1}{2} \end{pmatrix}\overrightharpoon{S(4)} + \begin{pmatrix} 1 \end{pmatrix}x(4) = \frac{1}{4} \\
	&\overrightharpoon{S(5)} = \begin{pmatrix} 0 & 1 \\ 0 & 0 \end{pmatrix}\overrightharpoon{S(4)} + \begin{pmatrix} 0 \\ 1 \end{pmatrix}x(4) = \overrightharpoon{0}& &y(5) = \begin{pmatrix} \frac{1}{4} & \frac{1}{2} \end{pmatrix}\overrightharpoon{S(5)} + \begin{pmatrix} 1 \end{pmatrix}x(5) = 0 \\
	&\vdots& &\vdots \\
\end{align*}

Thus, our system's outputs $y(n)$ for the input $x(n) = \delta(n) + \delta(n - 2)$ is given by:

\begin{align*}
	y(n) &= \delta(n) + \frac{1}{2}\delta(n - 1) + \frac{5}{4}\delta(n - 2) + \frac{1}{2}\delta(n - 3) + \frac{1}{4}\delta(n - 4) \\\\
	\intertext{\centering Or}
	y(n) &= \begin{cases} 1, \; \text{if} \; n = 0 \\ \frac{1}{2} \; \text{if} \; n = 1 \\ \frac{5}{4} \; \text{if} \; n = 2 \\ \frac{1}{2} \; \text{if} \; n = 3 \\ \frac{1}{4} \; \text{if} \; n = 4 \\0, \; \text{otherwise} \end{cases}
\end{align*}

\subsubsection[~\thesubsubsection]{Using convolution}

\begin{align*}
	y(n) &= \sum_{k = -\infty}^{k = \infty}(h(n - k)x(k)) \\
	&= \sum_{k = -\infty}^{k = \infty}((\delta(n - k) + \frac{1}{2}\delta(n - k - 1) + \frac{1}{4}\delta(n - k - 2))(\delta(k) + \delta(k - 2))) \\
	&= \sum_{k = -\infty}^{k = \infty}(\delta(n - k)\delta(k) + \frac{1}{2}\delta(n - k - 1)\delta(k) + \frac{1}{4}\delta(n - k - 2) + \delta(n - k)\delta(k - 2)
	\frac{1}{2}\delta(n - k - 1)\delta(k - 2) + \frac{1}{4}\delta(n - k - 2)\delta(k - 2)) \\
	&= \delta(0)\delta(n) + \frac{1}{2}\delta(0)\delta(n - 1) + \frac{1}{4}\delta(0)\delta(n - 2) + \delta(0)\delta{n - 2} + \frac{1}{2}\delta(0)\delta(n - 3) + \frac{1}{4}\delta(0)\delta(n - 4) \\
	&= \delta(n) + \frac{1}{2}\delta(n - 1) + \frac{5}{4}\delta(n - 2) + \frac{1}{2}\delta(n - 3) + \frac{1}{4}\delta(n - 4) \\
\end{align*}

Thus, our system's outputs $y(n)$ for the input $x(n) = \delta(n) + \delta(n - 2)$ is given by:

\begin{align*}
	y(n) &= \delta(n) + \frac{1}{2}\delta(n - 1) + \frac{5}{4}\delta(n - 2) + \frac{1}{2}\delta(n - 3) + \frac{1}{4}\delta(n - 4) \\\\
	\intertext{\centering Or}
	y(n) &= \begin{cases} 1, \; \text{if} \; n = 0 \\ \frac{1}{2} \; \text{if} \; n = 1 \\ \frac{5}{4} \; \text{if} \; n = 2 \\ \frac{1}{2} \; \text{if} \; n = 3 \\ \frac{1}{4} \; \text{if} \; n = 4 \\0, \; \text{otherwise} \end{cases}
\end{align*}

\newpage

\begin{center}
	\Large{\underline{Find the System}}
\end{center}

\section[Question ~\thesection]{For the system $y(n) = \alpha_0x(n) + \alpha_1x(n - 1) + \beta_1y(n - 1)$:}

\subsection[~\thesubsection]{Determine $\alpha_0, alpha_1$, and $\beta_1$ such that the impulse response $h(n) = \delta(n - 1)$}

\begin{align*}
	\alpha_0 &= 0 \\
	\alpha_1 &= 1 \\
	\beta_1 &= 0
\end{align*}

\subsection[~\thesubsection]{Determine $\alpha_0, alpha_1$, and $\beta_1$ such that the impulse response $h(n) = \frac{1}{3}^n \; n \ge 0$}

\begin{align*}
	\alpha_0 &= 1 \\
	\alpha_1 &= 0 \\
	\beta_1 &= \frac{1}{3}
\end{align*}

\subsection[~\thesubsection]{Determine $\alpha_0, alpha_1$, and $\beta_1$ such that the impulse response $h(n) = 1$}

\begin{align*}
	\alpha_0 &= 1 \\
	\alpha_1 &= 0 \\
	\beta_1 &= 1
\end{align*}

\newpage

\begin{center}
	\Large{\underline{Old Midterm Questions}}
\end{center}

\section[Question ~\thesection]{For the discrete signal $x(n) = \cos{(\frac{3\pi}{4}n)}$:}

\subsection[~\thesubsection]{What is the period of $x(n)$?}

The for a discrete periodic function, its period is defined as $P \in \mathbb{N}$ such that $x(n) = x(n + P) \; \forall n \in \mathbb{N}$. Using this fact, we find:

\begin{align*}
	x(n) &= x(n + P) \implies \\
	\cos{(\frac{3\pi}{4}n)} &= \cos{(\frac{3\pi}{4}(n + P))} \\
	&= \cos{(\frac{3\pi}{4}n + \frac{3\pi}{4}P)}
\end{align*}

Therefore, since cosine is $2\pi$ periodic, the smallest value of $P$ which makes $\frac{3\pi}{4}P$ a multiple of $2\pi$ is 8. Therefore, the period of $x(n)$ is 8.

\subsection[~\thesubsection]{What is the frequency of $x(n)$:}

The frequency of a cosine signal in the form $\cos{({\alpha}n)}$ is $\alpha$. Therefore, the frequency of $x(n)$ is $\frac{3\pi}{4}$.

\subsection[~\thesubsection]{What is the normalized frequency of $x(n)$?}

Here, the fundamental frequency of a discrete signal is the same as its frequency. Thus, the normalized frequency of $x(n)$ is $\frac{3\pi}{4}$.

\section[Question ~\thesection]{For the signal $x(n) = \sin{(\frac{2\pi}{3}n + \frac{\pi}{3})}$:}

\subsection[~\thesubsection]{What is the period of $x(n)$?}

Using the same method as in a.):

\begin{align*}
	x(n) &= x(n + P) \implies \\
	\sin{(\frac{2\pi}{3}n + \frac{\pi}{3})} &= \sin{(\frac{2\pi}{3}(n + P) + \frac{\pi}{3})} \\
	&= \sin{(\frac{2\pi}{3}n + \frac{2\pi}{3}P + \frac{\pi}{3})} \\\\
\end{align*}

Therefore, since cosine is $2\pi$ periodic, the smallest value of $P$ which makes $\frac{2\pi}{3}P$ a multiple of $2\pi$ is 3. Therefore, the period of $x(n)$ is 3.

\subsection[~\thesubsection]{What is the discrete frequency of $x(n)$?}

The discrete frequency of a signal is given by $f = \frac{1}{P}$ and thus the discrete frequency of $x(n)$ is $\frac{1}{3}$

\section[Question ~\thesection]{Given the system $y(n) = \frac{1}{2}(x(n) + x(n - 2))$, compute the output of the system for the signal $x(n) = 1 + \cos{({\pi}n)} + 3\sin{(\frac{\pi}{2}n + \pi)}$}

We can find the system's output by first finding the system's frequency response, $H(\omega)$, and use the formula $y(n) = |H(\omega)|\cos{({\omega}n + arg(H(\omega)))}$ (derived in lecture 7) to find the output of the system for each term in the input, and then summing them together (since the system has no feedback and is thus necessarily LTI)

\begin{directsubsubs}
\subsubsection[~\thesubsubsection]{Find the system's frequency response ($H(\omega)$)}

\begin{align*}
	y(n) &= \frac{1}{2}(x(n) + x(n - 2)) \\
	&= \frac{1}{2}x(n) + \frac{1}{2}x(n - 2) \implies \\
	H(\omega)e^{i{\omega}n} &= \frac{1}{2}e^{i{\omega}n} + \frac{1}{2}e^{i{\omega}(n - 2)} \\
	H(\omega)e^{i{\omega}n} &= \frac{1}{2}e^{i{\omega}n} + \frac{1}{2}e^{i{\omega}n - 2i{\omega})} \\
	H(\omega)e^{i{\omega}n} &= \frac{1}{2}e^{i{\omega}n} + \frac{1}{2}e^{i{\omega}n}e^{-2i{\omega})} \\
	H(\omega) &= \frac{1}{2} + \frac{1}{2}e^{-2i{\omega}} \\\\
\end{align*}

Therefore, our system's frequency response is given by $H(\omega) = \frac{1}{2} + \frac{1}{2}e^{-2i{\omega}}$.

\subsubsection[~\thesubsubsection]{Find the angular frequency of each term in $x(n)$}

%TODO: Reword

The angular frequency of all constant terms is 0 and the angular frequency of all sinusodials is given by the constant term beside n. Therefore:

\begin{tabularx}{\linewidth}{| c || Y |}
	\hline
	Terms of $x(n)$ & Angular Frequency $(\omega)$ \\
	\hline
	$1$ & $0$ \\
	$\cos{({\pi}n)}$ & $\pi$ \\
	$\sin{(\tfrac{\pi}{2}n + \pi)}$ & $\tfrac{\pi}{2}$ \\
	\hline
\end{tabularx}

\subsubsection[~\thesubsubsection]{Find $|H(\omega)|$ and $arg(H(\omega))$ for each term of $x(n)$}

\begin{tabularx}{\linewidth}{| c || Y | Y | Y | Y |}
	\hline
	Terms of $x(n)$ & Angular Frequency $(\omega)$ & Frequency Response $(H(\omega))$ & Gain $(|H(\omega)|)$ & Phase Shift $(arg(H(\omega)))$ \\
	\hline
	$1$ & $0$ & $1$ & $1$ & $0$ \\
	$\cos{({\pi}n)}$ & $\pi$ & $1$ & $1$ & $0$ \\
	$\sin{(\tfrac{\pi}{2}n + \pi)}$ & $\tfrac{\pi}{2}$ & $0$ & $0$ & - \\
	\hline
\end{tabularx}

\subsubsection[~\thesubsubsection]{Compute $y(n)$ using the aforementioned formula}

Note here that since $arg(H(\frac{\pi}{2}))$ does not exist, we don't include it in our computation.

\begin{align*}
	y(n) &= |H(\omega_1)|\cos{({\omega_1}n + arg(H(\omega_1))} + |H(\omega_2)|\cos{({\omega_2}n + arg(H(\omega_2))} \\
	&= \cos{((0)n + 0)} + \cos{((\pi)n + 0)} \\
	&= 1 + \cos{({\pi}n)} \\\\
\end{align*}

Therefore, the output to the system $y(n) = \frac{1}{2}(x(n) + x(n - 2))$ for the input $x(n) = 1 + \cos{({\pi}n)} + 3\sin{(\frac{\pi}{2}n + \pi)}$ is $1 + \cos{({\pi}n)}$

\end{directsubsubs}

\end{document}
